% Part: first-order-logic 
% Chapter: axiomatic-deduction 
% Section: rules-and-proofs

\documentclass[../../../include/open-logic-section]{subfiles}

\begin{document}

\iftag{FOL}
      {\olfileid{fol}{axd}{rul}}
      {\olfileid{pl}{axd}{rul}}

\olsection{قاعدے تے \usetoken{P}{derivation}}

\begin{explain}
  مسلّمی !!{derivation}s شاید منطق لئی سب توں سادہ !!{derivation} نظام
  نیں۔ !!^a{derivation} بس !!{formula}s دا اک سلسلہ اے۔ کسے سلسلے نوں
  !!a{derivation} منّن لئی اوہدے ہر !!{formula} دا یا تاں کسے مسلّمے
  دی اک مثال ہونا لازم اے، یا اوہدے توں پہلاں سلسلے وچ آ چکے اک یا ودھ
  !!{formula}s توں کسے استنباطی قاعدے نال نکلنا لازم اے۔
  !!^a{derivation} اپنے آخری !!{formula} نوں !!{derive}s کردا اے۔
\end{explain}

\begin{defn}[!!^{derivability}]
جے $\Gamma$، $\Lang L$ دے !!{formula}s دا سیٹ اے، تاں~$\Gamma$ توں
اک \emph{!!{derivation}}، !!{formula}s دا اک محدود سلسلہ $!A_1$،
\dots،~$!A_n$ اے، جیہدے ہر $i \le n$ لئی ہِٹھ دتیاں وچوں کوئی اک
گل درست ہووے:
\begin{enumerate}
\item $!A_i \in \Gamma$؛ یا
\item $!A_i$ اک مسلّمہ اے؛ یا
\item $!A_i$، کسے استنباطی قاعدے نال، کسے $!A_j$ (تے $!A_k$) توں
  نکلدا اے جتھے $j < i$ (تے $k < i$)۔
\end{enumerate}
\end{defn}

درست !!{derivation} کِنوں منّنا اے، ایہہ اس گل اُتے منحصر اے کہ اسیں
کیہڑے استنباطی قاعدے منّ دے آں (تے ظاہر اے کیہڑیاں گلاں نوں مسلّمے
منّ دے آں)۔ استنباطی قاعدہ اک جے-تاں بیان اے جیہڑا دسدَا اے کہ کِنہاں
شرطاں ہِٹھ !!a{derivation} دا قدم~$A_i$ اک درست استنباطی قدم اے۔

\begin{defn}[استنباط دا قاعدہ]
\emph{استنباط دا قاعدہ}،~$\Gamma$ توں !!a{derivation} وچ درست
استنباطی قدم ہون لئی اک کافی شرط دیندا اے۔
\end{defn}

مثال لئی، صرف اک عنصر دا سلسلہ $!A$، جتھے $!A \in \Gamma$، ظاہر اے
!!a{derivation} منّیا جاندا اے، ایس لئی ہِٹھ دتا بیان اک بہت سادہ
استنباطی قاعدہ ہو سکدا اے:
\begin{quote}
  جے $!A \in \Gamma$، تاں $!A$،~$\Gamma$ توں کسے وی
  !!{derivation} وچ ہمیشہ اک درست استنباطی قدم اے۔
\end{quote}
ایسے طرح جے $!A$ مسلّماں وچوں اک اے، تاں اکیلا $!A$ وی
!!a{derivation} اے؛ ایس لئی ایہ وی استنباطی قاعدہ اے:
\begin{quote}
  جے $!A$ اک مسلّمہ اے، تاں $!A$ اک درست استنباطی قدم اے۔
\end{quote}
قاعدہ ہور دلچسپ اُدوں ہوندا اے جدوں اوہ زیرِ غور قدم توں پہلاں آ چکے
!!{formula}s دا حوالہ دیندا اے۔ ہِٹھ دتے قاعدے نوں \emph{قاعدۂ اثباتِ
مقدم (modus ponens)} آکھدے نیں:
\begin{quote}
  جے $!B \lif !A$ تے $!B$، !!{derivation} وچ اِس توں اوپر آ چکے ہون،
  تاں~$!A$ اک درست استنباطی قدم اے۔
\end{quote}
جے ایہو اکّو استنباطی قاعدہ ہووے، تاں اُتّے دتی !!{derivation} دی
تعریف دا مطلب ایہ اے: $!A_1$، \dots،~$!A_n$ اُدوں تے صرف اُدوں
!!a{derivation} اے جدوں ہر $i \le n$ لئی ہِٹھ دتیاں وچوں کوئی اک
گل درست ہووے:
\begin{enumerate}
\item $!A_i \in \Gamma$؛ یا
\item $!A_i$ اک مسلّمہ اے؛ یا
\item کسے $j < i$ لئی $!A_j$، $!B \lif !A_i$ ہووے، تے کسے $k < i$
  لئی $!A_k$،~$!B$ ہووے۔
\end{enumerate}
آخری شق آکھدی اے کہ $!A_i$، قاعدۂ اثباتِ مقدم نال~$!A_j$
($!B \lif !A_i$) تے $!A_k$ ($!B$) توں نکلدا اے۔ جے اسیں $1$
توں~$n$ تک جا سکیے تے ہر واری ساڈے کول !!a{formula}~$!A_i$
ایہو جہا ہووے جیہڑا یا تاں~$\Gamma$ وچ ہووے، یا مسلّمہ ہووے، یا
جیہنوں کوئی استنباطی قاعدہ درست قدم دسے، تاں پورا سلسلہ اک درست
!!{derivation} منّیا جاندا اے۔

\begin{defn}[!!^{derivability}]
!!{formula}~$!A$، $\Gamma$ توں \emph{!!{derivable}} اے، جیہنوں
$\Gamma \Proves !A$ لکھدے آں، جے $\Gamma$ توں !!a{derivation}
موجود ہووے جیہڑا $!A$ اُتے ختم ہووے۔
\end{defn}

\begin{defn}[قضیے]
!!^a{formula}~$!A$ اک \emph{قضیہ} اے جے خالی سیٹ توں $!A$ دا
!!a{derivation} موجود ہووے۔ اسیں $\Proves !A$ لکھدے آں جے $!A$
قضیہ اے، تے جے نہیں تاں $\Proves/ !A$ لکھدے آں۔
\end{defn}


\end{document}
